\section{Conclusioni}\label{sec:conclusions}
La sensibilità al passo individua una regione favorevole per la ricetta di riferimento. I confronti fra architetture mostrano che l'esecuzione sincrona può ridurre il tempo conservando la traiettoria e che l'entità del guadagno dipende dal costo delle valutazioni. L'asincronia modifica la sequenza degli aggiornamenti e produce effetti diversi a seconda del problema.

Nello Studio III, il CPSO sincrono Random-8 con casualità per coordinata e $h=1$ fisso si colloca fra il terzo e il quarto posto su 13 metodi nel campione BBOB. A $50\,000$ FE ottiene un gap migliore del controllo discreto su 12 funzioni su 14 in ciascuna dimensione, con evidenza descrittiva dopo correzione multipla. CSO e SL-PSO mantengono un vantaggio in numerosi confronti. A tempo fissato, il multicore abbrevia il tempo al target nella diffusione inversa a $D=20$. Su Rastrigin, la ricerca raggiunge un plateau e l'aumento del numero di valutazioni incide poco sul gap finale.
